\section{Introducción}

El desarrollo de tecnologías de vanguardia exige una correspondencia estricta
entre el rigor analítico de la física teórica y la formulación de políticas de
Estado. Sin embargo, el campo de la nanotecnología exhibe una profunda fractura
entre su naturaleza física y su implementación institucional. Este
documento examina una cadena de errores categoriales que distorsionan el dominio
nanométrico, iniciando con la adopción de un mito historiográfico y culminando
en el colapso estructural de ecosistemas científicos en desarrollo. Se
argumenta que la imposición de métricas geométricas arbitrarias, en sustitución
de los criterios dictados por la mecánica cuántica, facilita una aproximación
puramente utilitaria de la ciencia. Al priorizar la aplicación
inmediata sobre el estudio de los principios fundamentales de la materia, las
políticas públicas generan un entorno de dependencia. El análisis se
centra en evidenciar cómo la suplantación del rigor físico por narrativas de
ingeniería macroscópica bloquea el progreso científico genuino, tomando como
paradigma empírico de este fracaso institucional el caso colombiano.
